% !TeX root = ../complete.tex
\section{مقایسه انرژی و تأخیر یک عملیات}
\label{sec:power-energy-delay}

\subsection{توان متوسط و انرژی به‌ازای هر عملیات}
\label{sec:power-energy-per-operation}

برای مقایسه دو پیاده‌سازی، ابتدا کاری را که هر دو باید انجام دهند
مشخص می‌کنیم. اگر مدار در بازه‌ای به مدت \(D\) فقط یک عملیات
انجام دهد و توان متوسط در همان مدت \(P_{\mathrm{avg}}\) باشد، از رابطه
\eqref{eq:power-interval}، انرژی عملیات \(E_{\mathrm{op}}=P_{\mathrm{avg}}D\)
به دست می‌آید. اگر توان را در بازه‌ای با چند عملیات اندازه بگیریم،
انرژی کل را بر تعداد عملیات کامل‌شده تقسیم می‌کنیم. با نرخ انجام
عملیات \(r_{\mathrm{op}}\)، یعنی تعداد عملیات کامل‌شده در ثانیه،
انرژی متوسط به‌ازای هر عملیات \(P_{\mathrm{avg}}/r_{\mathrm{op}}\)
است. در این کسر، واحد وات تقسیم بر معکوس ثانیه، ژول می‌دهد.
وقتی میانگین یک نتیجه در هر دوره ساعت باشد،
\(r_{\mathrm{op}}=f\) است. برای کاری که به چند دوره نیاز دارد،
رابطه را با نرخ واقعی انجام کار به کار می‌بریم
\cite[بخش‌های \secnum{۵.۴.۱} و \secnum{۵.۶}، صفحات
۲۰۰ و ۲۰۶]{weste}.

برای برآورد انرژی یک عملیات در تغذیه \(V\)، فرض کنید گره \(i\)
با ظرفیت معادل \(C_i\) در آن عملیات \(N_{\uparrow,i}\) بار از صفر
تا \(V\) شارژ می‌شود. مطابق رابطه \eqref{eq:power-source-energy}،
با جمع انرژی شارژِ همه گره‌ها، انرژی سوئیچینگ به دست می‌آید. مجموع
وزن‌دار \(C_{\mathrm{op}}=\sum_i N_{\uparrow,i} C_i\) واحد ظرفیت
دارد. اگر جریان نشتیِ متوسط \(I_{\mathrm{leak}}\) را برای مدت
انجام کار در نظر بگیریم و سهم پالس جریان مستقیم و هدایت هم‌زمان
ماندگار ناچیز باشند، مدل انرژی چنین است
\cite[رابطه \secnum{۵.۲۵}، صفحه ۲۰۲]{weste}.
\begin{equation}
  E_{\mathrm{op}}\approx E_{\mathrm{sw}}+E_{\mathrm{leak}}
    =C_{\mathrm{op}}V^2+I_{\mathrm{leak}}VD.
  \label{eq:power-operation-energy}
\end{equation}
جمله نخست با شمار شارژها و مربع ولتاژ تغییر می‌کند. در جمله دوم،
\(I_{\mathrm{leak}}V\) توان نشتی است و ضرب آن در مدت انجام کار،
انرژی نشتی را می‌دهد. بنابراین حتی کاهش توان نشتی می‌تواند با
بلندتر شدن زمان انجام کار به افزایش انرژی نشتی بینجامد.

\subsection{کمینه انرژی و کار در ولتاژ پایین}
\label{sec:power-minimum-energy}

در بخش \ref{sec:power-header-delay} از مدل جریان با اضافه‌ولتاژ
مثبت استفاده کردیم. اگر \(V\) پایین‌تر از \(V_t\) باشد، هنوز
جریان زیرآستانه وجود دارد. در تقریب بخش \secnum{۴.۴.۶.۴} کتاب،
جریان راه‌انداز با \(10^{V/S}\) متناسب است. \(S\) همان ولتاژ لازم
برای تغییر ده‌برابری جریان در مدل زیرآستانه است. با ثابت گرفتن
پارامترهای فناوری، ظرفیت مسیر و عوامل دیگری که بر جریان اثر دارند،
تقسیم بار متناسب با \(V\) بر جریان، تخمین زیر را برای مدت انجام
عملیاتِ محدود به یک مسیر می‌دهد
\cite[رابطه‌های \secnum{۴.۳۱} و \secnum{۵.۲۴}،
صفحات ۱۶۲ و ۲۰۲]{weste}.
\begin{equation}
  D(V)\approx K_s\,V\,10^{-V/S}.
  \label{eq:power-subthreshold-delay}
\end{equation}
\(K_s\) یک ضریب مثبت با واحد ثانیه بر ولت است که ظرفیت و ضریب
جریان مسیر را در بر می‌گیرد. اگر \(V\) به‌اندازه \(S\) کاهش
یابد، عامل نمایی ده برابر می‌شود. هم‌زمان، عامل خطی \(V\) کمی
کوچک‌تر می‌شود. در محدوده‌ای که اثر نمایی غالب است، تأخیر به‌سرعت
زیاد می‌شود. این مدل برای مقایسه نقطه‌های کاری معتبر در محدوده
زیرآستانه است. در ولتاژهای بسیار پایین، توانایی مدار در تولید و
حفظ سطوح منطقی نیز باید با معیار فصل ۲ بررسی شود.

برای نمونه، سه ولتاژ کاری زیر آستانه در یک مدل فرضی را مقایسه کنید.
می‌گیریم \(C_{\mathrm{op}}=\qty{20}{\femto\farad}\)،
\(I_{\mathrm{leak}}=\qty{5}{\nano\ampere}\) و
\(S=\qty{0.10}{\volt}\). برای این بازه فرض می‌کنیم تغییر
جریان نشتی با ولتاژ ناچیز است. زمان مرجع \(D_0=\qty{300}{\nano\second}\)
در \(V_0=\qty{0.30}{\volt}\) داده شده و همه‌ی ولتاژهای جدول در
محدوده معتبر مدل فرض می‌شوند. برای انتقال از زمان مرجع به زمان هر
ردیف، رابطه \eqref{eq:power-subthreshold-delay} را در دو ولتاژ بر
هم تقسیم می‌کنیم. حاصل به صورت زیر است.
\[
  \frac{D(V)}{D_0}\approx
    \frac{V}{V_0}\,10^{(V_0-V)/S}.
\]
در \qty{0.20}{\volt}، عامل نمایی ۱۰ و عامل خطی \(2/3\) است.
پس تأخیر تقریباً \((20/3)D_0=\qty{2000}{\nano\second}\) می‌شود.
در \qty{0.40}{\volt}، عامل نمایی \(1/10\) و عامل خطی \(4/3\)
است و تأخیر به \qty{40}{\nano\second} می‌رسد. برای نمونه در
ولتاژ پایین‌تر، جای‌گذاری در رابطه \eqref{eq:power-operation-energy}
دو مؤلفه را چنین می‌دهد.
\[
  \begin{aligned}
    E_{\mathrm{sw}}&=(\qty{20}{\femto\farad})(\qty{0.20}{\volt})^2
                     =\qty{0.80}{\femto\joule},\\
    E_{\mathrm{leak}}&=(\qty{5}{\nano\ampere})(\qty{0.20}{\volt})
                         (\qty{2.0}{\micro\second})
                     =\qty{2.0}{\femto\joule}.
  \end{aligned}
\]
همین محاسبه برای هر سه ولتاژ در جدول \ref{tab:power-energy-candidates}
آمده است. انرژی‌ها تا رقم اعشاری دوم نمایش داده شده‌اند تا جمع و
مقایسه ردیف‌ها روشن باشد.

\begin{table}[htbp]
  \centering
  \caption{انرژی و زمان یک عملیات در سه نقطه کاریِ یک مدل فرضی.
    \(E_{\mathrm{op}}\) از جمع انرژی سوئیچینگ و نشتی به دست آمده است.}
  \label{tab:power-energy-candidates}
  \begin{tabular}{ccccc}
    \toprule
    \(V\,(\unit{\volt})\) & \(D\,(\unit{\nano\second})\) &
    \(E_{\mathrm{sw}}\,(\unit{\femto\joule})\) &
    \(E_{\mathrm{leak}}\,(\unit{\femto\joule})\) &
    \(E_{\mathrm{op}}\,(\unit{\femto\joule})\) \\
    \midrule
    \num{0.20} & \num{2000} & \num{0.80} & \num{2.00} & \num{2.80} \\
    \num{0.30} & \num{300}  & \num{1.80} & \num{0.45} & \num{2.25} \\
    \num{0.40} & \num{40}   & \num{3.20} & \num{0.08} & \num{3.28} \\
    \bottomrule
  \end{tabular}
\end{table}

از میان این سه نقطه، \qty{0.30}{\volt} کم‌ترین انرژی را دارد.
در ولتاژ پایین‌تر، کاهش انرژی سوئیچینگ با افزایش انرژی نشتی در
زمان اجرای طولانی خنثی می‌شود. در ولتاژ بالاتر، انرژی سوئیچینگ
بزرگ‌تر است. این مقایسه کمینه را فقط در مجموعه سه گزینه مشخص
می‌کند. در مدار واقعی، تغییر فعالیت، آستانه یا دما می‌تواند
نقطه کمینه را جابه‌جا کند. این وابستگی در نقشه‌های شکل
\secnum{۵.۲۷} کتاب نیز دیده می‌شود \cite[صفحات ۲۰۲ و ۲۰۳]{weste}.

\subsection{کمینه \lr{Energy--Delay Product}}
\label{sec:power-edp}

گاهی هر دو هزینه انرژی عملیات و زمان اجرای آن در رتبه‌بندی اهمیت
دارند. معیار \lr{Energy--Delay Product} یا \lr{EDP}، حاصل‌ضرب
این دو با واحد ژول‌ثانیه است. محاسبه را برای \(V>V_t>0\) و با
\(C_{\mathrm{op}}>0\)، \(V_t\) و پارامترهای مسیر ثابت انجام می‌دهیم.
اگر انرژی سوئیچینگ غالب باشد، با ضرب انرژیِ رابطه
\eqref{eq:power-operation-energy} در تأخیرِ رابطه
\eqref{eq:power-above-threshold-delay} داریم
\cite[رابطه \secnum{۵.۲۶}، صفحه ۲۰۳]{weste}.
\begin{equation}
  \mathrm{EDP}(V)=E_{\mathrm{op}}(V)D(V)
       \approx K C_{\mathrm{op}}\frac{V^3}{(V-V_t)^m}.
  \label{eq:power-edp}
\end{equation}
برای یافتن کمینه، تغییر \(V^3(V-V_t)^{-m}\) را با ولتاژ بررسی
می‌کنیم. قاعده‌ی ضرب یک جمله از مشتق \(V^3\) و یک جمله از
مشتق \((V-V_t)^{-m}\) می‌دهد. با بیرون کشیدن عامل مشترک، مشتق
مدل \(\mathrm{EDP}\) چنین می‌شود.
\[
  \begin{aligned}
    \frac{\mathrm{d}(\mathrm{EDP})}{\mathrm{d}V}
       &=K C_{\mathrm{op}}\bigl[3V^2(V-V_t)^{-m}
                         -mV^3(V-V_t)^{-m-1}\bigr]\\
       &=K C_{\mathrm{op}} V^2 (V-V_t)^{-m-1}
           \bigl[(3-m)V-3V_t\bigr].
  \end{aligned}
\]
برای \(1\leq m\leq 2\)، عامل بیرون کروشه مثبت و شیب عبارت
خطی داخل کروشه \(3-m>0\) است. مشتق در نقطه‌ای که عبارت داخل
کروشه صفر می‌شود، از منفی به مثبت تغییر می‌کند. بنابراین کمینه
مدل در این ولتاژ است \cite[رابطه \secnum{۵.۲۷}، صفحه ۲۰۳]{weste}.
\begin{equation}
  V_{\mathrm{EDP}}=\frac{3}{3-m}\,V_t.
  \label{eq:power-edp-voltage}
\end{equation}
برای \(m=1.5\)، این ولتاژ دو برابر آستانه‌ است. اگر آستانه
فرضی \qty{0.30}{\volt} باشد، مدل \qty{0.60}{\volt} می‌دهد.
این محاسبه از تقریبِ صرف‌نظر از نشتی استفاده می‌کند. اگر ولتاژ
به گزینه‌های مشخص محدود باشد یا انرژی نشتی قابل‌توجه باشد، باید
انرژی و تأخیر را در گزینه‌های مجاز با مدل مناسب همان شرایط مقایسه کرد.

\subsection{کمینه انرژی با مهلت زمانی مشخص}
\label{sec:power-delay-constraint}

فرض کنید زمان انجام کار باید حداکثر \(D_{\max}\) باشد. ابتدا فقط
گزینه‌هایی را نگه می‌داریم که \(D\leq D_{\max}\) دارند و سپس انرژی
آن‌ها را مقایسه می‌کنیم. برای مثال، با مهلت \qty{100}{\nano\second}
از جدول \ref{tab:power-energy-candidates} فقط گزینه‌ی
\qty{0.40}{\volt} می‌ماند. اگر مهلت به \qty{300}{\nano\second}
برسد، گزینه‌ی \qty{0.30}{\volt} نیز مجاز می‌شود و میان گزینه‌های
مجاز انرژی کم‌تری دارد.

این روش برای مقایسه توپولوژی و اندازه گیت‌ها نیز کاربرد دارد.
در ولتاژ ثابت، از \NotesTopic{sec:delay-path-model}{مدل تأخیر مسیر}
برای یافتن کندترین مسیر هر گزینه و از مجموع ظرفیت‌های وزن‌دار
رابطه \eqref{eq:power-operation-energy} برای برآورد انرژی استفاده
می‌کنیم. همان‌طور که مثال بخش \ref{sec:power-capacitance-reduction}
نشان داد، کاهش اندازه یک گیت ممکن است انرژی را کم و تأخیر مسیر را
زیاد کند. کتاب شکل \secnum{۵.۱۳} را با تکرار همین مسئله برای
مهلت‌های متفاوت به دست می‌آورد. هر نقطه روی آن منحنی، انرژیِ کمینه
مدل با یک مهلت مشخص را نشان می‌دهد. بنابراین مقدار مهلت بخشی از
تعریف مسئله‌ی طراحی است \cite[بخش‌های \secnum{۵.۲.۲.۱} و
\secnum{۵.۴.۳}، صفحات ۱۸۹، ۱۹۰، ۲۰۳ و ۲۰۴]{weste}.
